% Part: lambda-calculus
% Chapter: representability
% Section: currying

\documentclass[../../../include/open-logic-section]{subfiles}

\begin{document}

\olfileid{lam}{rep}{cur} 
\olsection{কারিকরণ}

কোনো $\lambd$-বিমূর্তন $\lambd[x][M]$ একটিমাত্র আর্গুমেন্টের অপেক্ষককে
প্রকাশ করে; অথচ একাধিক আর্গুমেন্ট গ্রহণকারী অপেক্ষক সংজ্ঞায়িত করতে চাইলে
এটি বেশ বড় একটি সীমাবদ্ধতা। এর একটি উপায় হতে পারত $\lambd$-ক্যালকুলাসকে
সম্প্রসারিত করে জোড়, ত্রয়ী ইত্যাদি গঠনের অনুমতি দেওয়া; সে ক্ষেত্রে,
ধরা যাক, কোনো ত্রি-স্থানীয় অপেক্ষক $\lambd[x][M]$ তার আর্গুমেন্ট হিসেবে
একটি ত্রয়ী আশা করত। কিন্তু \emph{কারিকরণ}-এর সাহায্যে কাজটি করা আরও
সুবিধাজনক।

একটি উদাহরণ বিবেচনা করি। মুহূর্তের জন্য ধরে নিই, $\lambd$-ক্যালকুলাসে
একটি~$+$ ক্রিয়া আছে। যোগ অপেক্ষকটি $2$-স্থানীয়, অর্থাৎ সেটি দুটি
আর্গুমেন্ট নেয়। কিন্তু একটি $\lambd$-বিমূর্তন থেকে আমরা কেবল এক
আর্গুমেন্টের অপেক্ষক পাই: বাক্যগঠন $\lambd[(x,y)][(x+y)]$-এর মতো রাশির
অনুমতি দেয় না। তবে $\lambd[y][(x+y)]$ দ্বারা প্রদত্ত এক-স্থানীয়
অপেক্ষক~$f_x(y)$ বিবেচনা করা যায়, যা তার একমাত্র আর্গুমেন্ট~$y$-এর সঙ্গে
$x$ যোগ করে। আসলে এটি একটি অপেক্ষক নয়, বরং প্রতিটি সংখ্যা~$x$-এর জন্য
একটি করে ``$x$ যোগ করো'' অপেক্ষকের পরিবার। এবার আরেকটি এক-স্থানীয়
অপেক্ষক~$g$-কে $\lambd[x][f_x]$ হিসেবে সংজ্ঞায়িত করা যায়। আর্গুমেন্ট
$x$-এর উপর প্রয়োগ করলে $g(x)$ অপেক্ষক~$f_x$-কে ফেরত দেয়---অর্থাৎ তার
মানগুলিও অপেক্ষক। এখন $g$-কে~$x$-এর উপর এবং তারপর ফলটিকে~$y$-এর উপর
প্রয়োগ করলে পাই: $(g(x))y = f_x(y) = x+y$। এইভাবে এক-স্থানীয়
অপেক্ষক~$g$ দ্বি-স্থানীয় যোগ অপেক্ষকেরই কাজ করতে পারে। ``কারিকরণ'' বলতে
শুধু দ্বি-স্থানীয় অপেক্ষককে এমন এক-স্থানীয় অপেক্ষকে বদলানোর এই কৌশলই
বোঝায়, যার মানগুলি আবার এক-স্থানীয় অপেক্ষক।

এবার $\lambd$-ক্যালকুলাসের যথাযথ বাক্যগঠনে একটি উদাহরণ দেখা যাক। অপেক্ষক
$f(x,y) = x$-কে কীভাবে প্রকাশ করব? দুটি আর্গুমেন্ট গ্রহণ করে প্রথমটি
ফেরত দেয় এমন অপেক্ষক সংজ্ঞায়িত করতে আমরা
$\lambd[x][\lambd[y][x]]$ লিখতে পারি। আক্ষরিক অর্থে এটি এমন একটি অপেক্ষক,
যা একটি আর্গুমেন্ট~$x$ গ্রহণ করে এবং অপেক্ষক~$\lambd[y][x]$-কে ফেরত দেয়।
অপেক্ষক $\lambd[y][x]$ আরেকটি আর্গুমেন্ট~$y$ গ্রহণ করে, কিন্তু সেটিকে
উপেক্ষা করে সব সময়~$x$-ই ফেরত দেয়। $\lambd[x][\lambd[y][x]]$-কে দুটি
আর্গুমেন্টের উপর প্রয়োগ করলে কী হয় দেখা যাক:
\begin{align*}
  (\lambd[x][\lambd[y][x]])MN 
  \bredone & (\lambd[y][M])N \\
  \bredone & M
\end{align*}

সাধারণভাবে, কোনো পদ~$N$ দ্বারা সংজ্ঞায়িত $x_1$, \dots,~$x_n$ পরামিতির
অপেক্ষক লিখতে পারি
$\lambd[x_1][\lambd[x_2][\ldots\lambd[x_n][N]]]$। এর উপর $n$টি আর্গুমেন্ট
প্রয়োগ করলে পাই:
\begin{multline*}
  (\lambd[x_1][\lambd[x_2][\ldots\lambd[x_n][N]]]) M_1 \dots M_n \bredone\\
  \begin{aligned}
  \bredone {} & (\Subst{(\lambd[x_2][\ldots\lambd[x_n][N]])}{M_1}{x_1}) M_2
  \dots M_n\\
   \eqs {} & (\lambd[x_2][\ldots\lambd[x_n][\Subst{N}{M_1}{x_1}]]) M_2
            \dots M_n \\
  \vdots & \\
  \bredone {} &\Subst{\Subst{N}{M_1}{x_1}\ldots}{M_n}{x_n}
  \end{aligned}
\end{multline*}
% BN-SRC-267 TARGET-CORRECTION-BEGIN
\par\smallskip\noindent\textnormal{\small\textbf{উৎস-সংশোধন (BN-SRC-267):} হিমায়িত সূত্রের শেষ হ্রাসে অপেক্ষকের সংজ্ঞায় ব্যবহৃত দেহ~$N$-এর বদলে অনির্দিষ্ট~$P$ লেখা হয়েছে। বাংলা সূত্রে~$N$ পুনরুদ্ধার করা হয়েছে, যাতে ধারাবাহিক প্রতিস্থাপনগুলি ঠিক মূল অপেক্ষক-দেহেই প্রয়োগ হয়।} % BN-SRC-267 TARGET-CORRECTION-END
শেষ পঙ্‌ক্তিটির আক্ষরিক অর্থ হলো অপেক্ষক-সংজ্ঞার দেহে প্রতিটি~$x_i$-এর
স্থানে~$M_i$ প্রতিস্থাপন করা; একাধিক আর্গুমেন্ট কোনো অপেক্ষকের উপর প্রয়োগ
করলে আমরা ঠিক এটিই চাই।
\end{document}

